% Macros lifted from snake.tex's preamble, with the xy-pic and tikz-cd definitions
% dropped: plasTeX cannot parse them, and each diagram is instead compiled to SVG by
% real LaTeX (with the paper's own preamble, \newdir lines included) and included as
% an image.  This file is maintained by hand: a macro new to snake.tex has to be
% copied here, or the blueprint renders it undefined.  Definitions here must also be ones
% MathJax understands: \mathpzc, \xspace and lengths such as \horspace are not, and leak
% into the page as literal text.  MathJax has the AMS packages and nothing more, so a macro
% that can only expand to mathtools (\xRightarrow, \coloneq) keeps its mathtools spelling
% here, where the pdf version reads it, and is shimmed for the web in macros/web.tex.
% The paper numbers its results by subsection, sharing one counter; reproducing that
% here makes every blueprint number equal to the number printed in the paper.
\newtheorem{theorem}[subsection]{Theorem}
\newtheorem{proposition}[subsection]{Proposition}
\newtheorem{lemma}[subsection]{Lemma}
\newtheorem{corollary}[subsection]{Corollary}
\theoremstyle{definition}
\newtheorem{definition}[subsection]{Definition}
\newtheorem{example}[subsection]{Example}
\newtheorem{remark}[subsection]{Remark}
\newtheorem{convention}[subsection]{Convention}
\newtheorem{notation}[subsection]{Notation}
\newenvironment{tfae}{\begin{enumerate}}{\end{enumerate}}
\newenvironment{enum}{\begin{enumerate}}{\end{enumerate}}
\newenvironment{enumT}{\begin{enumerate}}{\end{enumerate}}
% `eqD` is rewritten to `equation` by genblueprint.py, which decides per display whether
% its body is a diagram image (then `center`) or a formula (then MathJax); no shim here.

\newcommand{\noproof}{\hfill \qed}
\def\nameit#1{\textrm{#1}~}
\def\thex{\nameit{Theorem}}
\def\prox{\nameit{Proposition}}
\def\corx{\nameit{Corollary}}
\def\lemx{\nameit{Lemma}}
\def\defx{\nameit{Definition}}
\def\remx{\nameit{Remark}}
\def\exax{\nameit{Example}}
\def\dfn#1{{\itshape #1}}
\newcommand{\refs}[1]{\textup{(}\ref{#1}\textup{)}}
\newcommand{\h}[1][1.0]{\,}
\def\:{\colon}
\def\c{\circ}
\newcommand{\comp}{\circ}
\newcommand{\iso}{\cong}
\newcommand{\ov}[1]{\overline{#1}}
\def\phi{\varphi}
\def\eps{\varepsilon}
\newcommand{\catfont}[1]{\ensuremath{\mathit{#1}}}
\newcommand{\onecat}[1]{\ensuremath{\mathsf{#1}}}
\def\L{\catfont{L}}
\newcommand{\Cat}{\catfont{Cat}}
\newcommand{\AbCat}{\catfont{AbCat}}
\newcommand{\Triang}{\catfont{Triang}}
\newcommand{\Sat}{\catfont{Sat}}
\newcommand{\LAdj}{\catfont{LAdj}}
\newcommand{\Sup}{\catfont{Sup}}
\newcommand{\Supmod}{\ensuremath{\mathit{Sup}_{\mathrm{mod}}}}
\newcommand{\Suphil}{\ensuremath{\mathit{Sup}_{\mathrm{hil}}}}
\renewcommand{\1}{\onecat{1}}
\newcommand{\Grp}{\onecat{Grp}}
\newcommand{\Ab}{\onecat{Ab}}
\newcommand{\Perf}{\onecat{Perf}}
\newcommand{\HypAb}{\onecat{HypAb}}
\newcommand{\Nil}{\onecat{Nil}}
\newcommand{\CMon}{\onecat{CMon}}
\newcommand{\SES}{\onecat{SES}}
\newcommand{\VectF}{\ensuremath{\mathsf{Vect}_{F}}}
\newcommand{\Mlc}{\onecat{Mlc}}
\newcommand{\Coh}{\onecat{Coh}}
\newcommand{\Scl}{\ensuremath{\mathrm{Sc}}}
\newcommand{\Lat}{\ensuremath{\mathrm{L}}}
\newcommand{\dseg}[1]{\ensuremath{{\downarrow}#1}}
\newcommand{\useg}[1]{\ensuremath{{\uparrow}#1}}
\newcommand{\trunc}{\ensuremath{\tau_{1}}}
\DeclareMathOperator{\Hom}{Hom}
\DeclareMathOperator{\End}{End}
\DeclareMathOperator{\Ext}{Ext}
\DeclareMathOperator{\Spec}{Spec}
\DeclareMathOperator{\Supp}{Supp}
\DeclareMathOperator{\Ass}{Ass}
\DeclareMathOperator{\ASpec}{ASpec}
\DeclareMathOperator{\ASupp}{ASupp}
\DeclareMathOperator{\AAss}{AAss}
\newcommand{\AS}{\textup{(AS)}}
\newcommand{\DPN}{\textup{(DPN)}}
\newcommand{\NEC}{\textup{(NEC)}}
\newcommand{\Ker}{\ensuremath{\mathrm{Ker}}}
\newcommand{\Coker}{\ensuremath{\mathrm{Cok}}}
\newcommand{\coker}{\ensuremath{\mathrm{coker}}}
\newcommand{\TwoKer}{\ensuremath{2\text{-}\mathrm{Ker}}}
\newcommand{\TwoCoker}{\ensuremath{2\text{-}\mathrm{Cok}}}
\newcommand{\twoker}{\ensuremath{2\text{-}\mathrm{ker}}}
\newcommand{\twocoker}{\ensuremath{2\text{-}\mathrm{coker}}}
\renewcommand{\Im}{\ensuremath{\mathrm{Im}}}
\newcommand{\im}{\ensuremath{\mathrm{im}}}
\newcommand{\Coim}{\ensuremath{\mathrm{Coim}}}
\newcommand{\coim}{\ensuremath{\mathrm{coim}}}
\newcommand{\TwoImg}{\ensuremath{2\text{-}\mathrm{Im}}}
\newcommand{\twoimg}{\ensuremath{2\text{-}\mathrm{im}}}
\newcommand{\TwoCoim}{\ensuremath{2\text{-}\mathrm{Coim}}}
\newcommand{\twocoim}{\ensuremath{2\text{-}\mathrm{coim}}}
\newcommand{\Hker}[2]{\ensuremath{\mathrm{H}^{\ker}_{#1}(#2)}}
\newcommand{\Hcoker}[2]{\ensuremath{\mathrm{H}^{\coker}_{#1}(#2)}}
\newcommand{\HomC}[3]{{#1}\left({#2},\h[1]{#3}\right)}
\newcommand{\id}[1]{\operatorname{id}_{#1}}
\newcommand{\op}{\ensuremath{^{\operatorname{op}}}}
\newcommand{\too}{\longrightarrow}
\newcommand{\ito}{\hookrightarrow}
\newcommand{\iito}{\rightarrowtail}
\newcommand{\toe}{\twoheadrightarrow}
\newcommand{\aar}[2][]{\xrightarrow[#1]{#2}}
% The paper builds these two extensible arrows from amsmath internals (\ext@arrow,
% \Rightarrowfill@), which MathJax cannot expand and prints as red literal text.  These
% are the same arrows from MathJax's extpfeil package, which its autoload provides.
\newcommand{\aR}[2][]{\xRightarrow[#1]{#2}}
\newcommand{\aeqq}[2][]{\xlongequal[#1]{#2}}
\newcommand*{\comment}[4][]{%
  \@ifundefined{c@#2}{\newcounter{#2}}{}%
  \begingroup
  \ifblank{#1}{%
    \newcommand*{\header}{\csname the#2\endcsname}%
  }{%
    \newcommand*{\header}{(#1\csname the#2\endcsname)}%
  }%
  \addtocounter{#2}{1}%
  \textnormal{\textcolor{#3}{\header.[#4\@]}}%
  \endgroup
}
\newcommand*{\elena}[1]{\comment[E]{elena}{orange}{#1}}
\newcommand*{\tim}[1]{\comment[T]{tim}{teal}{#1}}
\newcommand*{\luca}[1]{\comment[L]{luca}{brown}{#1}}
\newcommand*{\todo}[1]{\comment{todo}{blue}{#1}}
